% Complete actual-class completion, 19 September 2026.
% Original DNE and FCI sources are unchanged.
\section{The full arithmetic class, its joint Gram, and its action}

Retain the simple quartet, the original period and branch, the full
nonzero conductor $E_A$, and the original full arithmetic unit.
Write $a=k+4$, $k=4\ell+1\ge9$,
$q=(a+1)^2$, $q'=(a-7)^2$, $v=\operatorname{ord}_0E_A\le80$,
$m=q-q'-v$, $c=a/2$, and
\[
 \beta_{a,r,t}=c+(2r-a)\delta+i(2t-a)\gamma,\qquad
 0<\delta<\tfrac12,\quad \gamma>2,\qquad
 \chi(S)=\prod_{r,t=0}^{a}(S-\beta_{a,r,t}).
 \tag{NCRP1}
\]
In particular $q'>0$, $m>0$, every root is simple, and
$\Re\beta\ge a(1/2-\delta)>0$. No root or conductor factor is
removed. The residue-dual polynomial frame and original DNE minor
are those proved in the preceding section. Let
$\mathscr P:\mathbb C^q\to E=\mathbb C[S]/(\chi)$ send the standard
column $e_n$ to $\mathfrak p_n$.
Thus $\mathscr P=\mathcal L\mathbb V^{-1}$, exactly in FCI33--34,
and multiplication by $S$ has matrix $C$ in this frame:
\[
 C_{\ell n}=\mathbf1_{\{n\ge1,\ell=n-1\}}
              -c_n\mathbf1_{\{\ell=q-1\}},\qquad
 \chi(S)=\sum_{n=0}^{q}c_nS^n,\quad c_q=1.
 \tag{NCRP2}
\]
This use of $C$ denotes the full packet action matrix; the original
conductor coefficient matrix continues to be written $C_a$.

\subsection{A square completion using the actual minor}

Let $E_{<v}:\mathbb C^v\to\mathbb C^q$ insert the first $v$
coordinates and retain the original $P_v,J_v,\mathcal I,\mathcal J$.
Define the square coefficient map
\[
 K_{\mathcal I}=
 [E_{<v},\,P_vJ_v,\,P_vE_{\mathcal J}],\qquad
 F_{\mathcal I}=\mathscr P K_{\mathcal I}=[F_0,F_C,L_a].
 \tag{NCRP3}
\]
Here $F_0$ has $v$ columns, $F_C=\mathcal L C_a^{\mathsf T}$
has $q'$ columns, and $L_a=\mathcal L Z_{\exp}$ has $m$ columns.
All three maps have codomain the original polynomial packet $E$.
For a moment column $t=(t_{<v},t_{\ge v})$ the inverse is explicitly
\[
 K_{\mathcal I}^{-1}t=
 \begin{pmatrix}
 t_{<v}\\
 (J_v)_{\mathcal I}^{-1}(t_{\ge v})_{\mathcal I}\\
 (t_{\ge v})_{\mathcal J}
       -(J_v)_{\mathcal J}(J_v)_{\mathcal I}^{-1}
                                (t_{\ge v})_{\mathcal I}
 \end{pmatrix}.
 \tag{NCRP4}
\]
Indeed substitution in the $\mathcal I$ and $\mathcal J$ rows,
and then in the first $v$ rows, gives exactly $t$.
Thus $K_{\mathcal I}$ and $F_{\mathcal I}$ are isomorphisms.
This proves a full arithmetic representative map, not only a
relative quotient map of dimension $m$.

For a retained simple eigenclass $v_\lambda$ let
$t_{\lambda,n}=\lambda^n/\chi'(\lambda)$ and
\[
 w_\lambda=(u_0,\eta,z)^{\mathsf T}
                  =K_{\mathcal I}^{-1}t_\lambda .
 \tag{NCRP5}
\]
The preceding residue calculation gives the exact equality
$v_\lambda=F_0u_0+F_C\eta+L_az$.
The vector $\eta$ is nonzero, because $q'>0$, all entries of
$(t_{\lambda,\ge v})_{\mathcal I}$ are nonzero, and the retained
minor is invertible. If $v>0$, $u_0$ is also nonzero.
Since $F_{\mathcal I}$ is injective, the omitted vector
$F_0u_0+F_C\eta$ is nonzero. These assertions hold for the actual
minor selected in DNE; they do not replace it by the consecutive chart.

\subsection{The original full-unit arithmetic injection}

The map into theta cohomology retains the source $F_h$ of FCI35:
$\mathcal M F_h=v_h=2\xi/h$ and $h(D)F_h=\Theta\phi_*$.
On the simple-quartet branch,
$h(s)=\prod_{\alpha\in\{\rho,\bar\rho,1-\rho,1-\bar\rho\}}(s-\alpha)$.
For $P\in E$ set, with the original Euler variables,
\[
 \iota_a(P)=
 [r_\chi(P)(D_1+\cdots+D_a)F_h^{\otimes a}]
            \in Q_\theta^{\otimes a}.
 \tag{NCRP6}
\]
Here $r_\chi$ is the increasing-degree representative below $q$.
We recall the complete finite reason that this packet map is injective.
In the tensor algebra
$\mathbb C[s_1,\ldots,s_a]/(h(s_1),\ldots,h(s_a))$,
evaluation at all ordered root tuples is an isomorphism: apply the
four-root Vandermonde separately in every coordinate.
Every pair of counts $(r,t)$, $0\le r,t\le a$, occurs among such
tuples, since the real-sign and imaginary-sign choices can be made
independently in each tensor leg. The resulting distinct sums are
exactly the roots in NCRP1. Hence $P(s_1+\cdots+s_a)$ vanishes
in that tensor algebra precisely when $\chi$ divides $P$.
The cyclic annihilator is therefore the full $\chi$.

The original Taylor map on the theta quotient sends NCRP6 to
$\upsilon_h^{\otimes a}P(A_h^{(a)})1$.
Every component of the original primary unit $\upsilon_h=j_h(v_h)$
is nonzero, so multiplication by it is invertible. It cannot turn
a nonzero packet class into zero.
For completeness, multiples of $\chi(s_1+\cdots+s_a)$ lie in the
ideal $(h(s_1),\ldots,h(s_a))$ by the preceding tensor evaluation.
Successive monic division in $s_1,\ldots,s_a$ writes such a polynomial
as $\sum_i h(s_i)R_i$. Substitution of the original differential
operators supplies its theta primitive, leg by leg, with the original
$(-1)^{i-1}$ tensor-differential signs. Thus the representative
independence and the Taylor injectivity apply to the same map.
This is the FCI35/FDC construction \cite{prog:FCI}, with its full
unit and all ordered tuples retained.

Consequently
\[
 \iota_a(v_\lambda)
 =\iota_a(F_0u_0)+\iota_a(F_C\eta)+\iota_a(L_az),
 \qquad
 \iota_a(F_0u_0+F_C\eta)\ne0.
 \tag{NCRP7}
\]
Taking four independent copies of these maps preserves every
equation. The omitted summands cannot be removed as theta boundaries.

\subsection{The full original arithmetic metric}

Let $dm_a^{\rm ar}(y)=m_a^{\rm ar}(y)\,dy$,
$m_a^{\rm ar}=w_h^{*a}$, and
$w_h(y)=|v_h(1/2+iy)|^2/(2\pi)$.
The source has its complete mass $(\int_{\mathbb R}w_h)^a$.
Let $H_N^{\rm ar}$ be its moment Gram in the increasing powers of
$S=c+iy$. For the literal remainder map
$J_N^{\chi}:\mathbb C[S]_{\le N}\to E$ and $N\ge q-1$, set
\[
 G_N=(J_N^\chi(H_N^{\rm ar})^{-1}(J_N^\chi)^*)^{-1},
 \qquad
 T_N=(H_N^{\rm ar})^{-1}(J_N^\chi)^*G_N.
 \tag{NCRP8}
\]
The norm of a class is the minimum over its full relation fibre.
To verify the formula, $J_N^\chi T_N=I$ and
$(T_Nx)^*H_N^{\rm ar}r=0$ whenever $J_N^\chi r=0$.
Every other lift is $T_Nx+r$, so its squared norm is
$x^*G_Nx+r^*H_N^{\rm ar}r$. This proves both the minimum and
$T_N^*H_N^{\rm ar}T_N=G_N$ with the whole relation space
$\chi\mathbb C[S]_{\le N-q}$ retained. It is the classical positive
Schur calculation \cite{human:boyd-vandenberghe}, proved here in
its complex Hermitian coordinates.

All polynomial moments are finite for the admitted arithmetic source.
Its density is positive almost everywhere: the nonzero analytic
function $v_h$ on the physical line has isolated zeros; hence
$w_h>0$ almost everywhere, and its convolution powers are positive.
It follows that every nonzero polynomial has positive squared norm.
Thus all displayed moment and quotient forms are positive definite.
The original physical coordinate, both integration half-lines, and
the full root polynomial have remained in NCRP8.

Write $\mathbf F_{\mathcal I}$ for the increasing-power coefficient
matrix of $F_{\mathcal I}$. Its complete Gram and action are
\[
 \mathsf H_N=\mathbf F_{\mathcal I}^*G_N\mathbf F_{\mathcal I},
 \qquad
 \mathsf A=K_{\mathcal I}^{-1}CK_{\mathcal I}.
 \tag{NCRP9}
\]
The map $F_{\mathcal I}$ intertwines $\mathsf A$ and multiplication
by $S$ exactly, so
$\mathsf A w_\lambda=\lambda w_\lambda$.
Partition the blocks by $0,C,Z$, of sizes $v,q',m$. The actual
squared norm of the retained class is
\[
 \begin{split}
 \|v_\lambda\|_{G_N}^2={}&
 u_0^*\mathsf H_{00}u_0+\eta^*\mathsf H_{CC}\eta
       +z^*\mathsf H_{ZZ}z\\
 &+2\Re\{u_0^*\mathsf H_{0C}\eta
              +u_0^*\mathsf H_{0Z}z
              +\eta^*\mathsf H_{CZ}z\}.
 \end{split}                                                     \tag{NCRP10}
\]
These are the three diagonal terms and all six directed cross
pairings. There is no orthogonality claim about the DNE decomposition.

The full centered action has blocks
\[
 \mathsf D_{ij}
 =\sum_{r\in\{0,C,Z\}}\mathsf A_{ri}^*\mathsf H_{rj}
  +\sum_{r\in\{0,C,Z\}}\mathsf H_{ir}\mathsf A_{rj}
  -a\mathsf H_{ij}.
 \tag{NCRP11}
\]
Direct multiplication gives
$\mathsf D=\mathsf A^*\mathsf H+\mathsf H\mathsf A-a\mathsf H$.
For $\lambda=a\rho=c+a\delta+ia\gamma$, the original same-class
Hermitian action therefore obeys the exact full receiver
\[
 \begin{split}
 &u_0^*\mathsf D_{00}u_0+\eta^*\mathsf D_{CC}\eta
       +z^*\mathsf D_{ZZ}z\\
 &\qquad+
 2\Re\{u_0^*\mathsf D_{0C}\eta+
                  u_0^*\mathsf D_{0Z}z+
                  \eta^*\mathsf D_{CZ}z\}
       =2a\delta\, w_\lambda^*\mathsf H_Nw_\lambda>0.
 \end{split}                                                     \tag{NCRP12}
\]
The proof substitutes $\mathsf Aw_\lambda=\lambda w_\lambda$;
both terms contribute their full $\lambda,\bar\lambda$.
This is a statement on the complete retained class. It has not
assigned a sign to one of its summands.

\subsection{Orthogonal completion with every original cross block}

It is useful to obtain an exact minimum-distance formula without
changing the original minor. Put
$F_\perp^{\,0}=[F_0,F_C]$ and $u=(u_0,\eta)^{\mathsf T}$.
In the coefficient formulas below the maps are expressed in the
same increasing-power frame used for $G_N$. Define
\[
 \begin{gathered}
 \Gamma=L_a^*G_NL_a,\qquad
 H_{ZC}=L_a^*G_NF_\perp^{\,0},\qquad
 H_{CC}=(F_\perp^{\,0})^*G_NF_\perp^{\,0},\\
 P_Z=L_a\Gamma^{-1}L_a^*G_N,\qquad
 E_C=(I-P_Z)F_\perp^{\,0},\\
 \Sigma=H_{CC}-H_{ZC}^*\Gamma^{-1}H_{ZC},\qquad
 \widehat z=z+\Gamma^{-1}H_{ZC}u .
 \end{gathered}                                                   \tag{NCRP13}
\]
Completing the positive square proves that $\Sigma\succ0$:
for $u\ne0$, $E_Cu=0$ would place $F_\perp^{\,0}u$ in
$\operatorname{im}L_a$, contradicting the invertibility of
the full frame in NCRP3. Since $P_Z$ is the exact $G_N$-orthogonal
projector, $L_a^*G_NE_C=0$ and $E_C^*G_NE_C=\Sigma$.
Substitution now gives
\[
 \begin{gathered}
 v_\lambda=L_a\widehat z+E_Cu,\qquad
 \|v_\lambda\|_{G_N}^2
       =\widehat z^*\Gamma\widehat z+u^*\Sigma u,\\
 \operatorname{dist}_{G_N}(v_\lambda,\operatorname{im}L_a)^2
       =u^*\Sigma u>0.
 \end{gathered}                                                   \tag{NCRP14}
\]
This is a quantitative exact repair of the omitted-class issue.
The cross blocks in NCRP10 have contributed to both $\widehat z$
and $\Sigma$; they were not set to zero.

Write the action in the original algebraic frame, now ordered by
$(C,Z)$, as
\[
 A F_\perp^{\,0}=F_\perp^{\,0}A_{CC}+L_aA_{ZC},\qquad
 A L_a=F_\perp^{\,0}A_{CZ}+L_aA_{ZZ}.
 \tag{NCRP15}
\]
These matrices are the corresponding blocks of NCRP9 and therefore
are explicitly calculated from $K_{\mathcal I}^{-1}CK_{\mathcal I}$.
In the metric-orthogonal full frame $[L_a,E_C]$ they give
\[
 \begin{gathered}
 A_{\rm nat}=\Gamma^{-1}L_a^*G_NAL_a
              =A_{ZZ}+\Gamma^{-1}H_{ZC}A_{CZ},\\
 R_N=AL_a-L_aA_{\rm nat}=E_CA_{CZ},\\
 B_N=\Gamma^{-1}L_a^*G_NAE_C,\qquad
 D_N^{C}=\Sigma^{-1}E_C^*G_NAE_C,\\
 [A]_{[L_a,E_C]}
       =\begin{pmatrix}A_{\rm nat}&B_N\\ A_{CZ}&D_N^{C}\end{pmatrix}.
 \end{gathered}                                                   \tag{NCRP16}
\]
The $D_N^{C}$ here is a complement action block, not the native
Gamma metric named $D_N$ in earlier source modules.
For the lower-left block, subtracting $P_ZAL_a$ gives
$E_CA_{CZ}$, and multiplication by the inverse of $\Sigma$
proves its coefficient is exactly $A_{CZ}$.
Thus the full eigenclass has the two coupled equations
\[
 (A_{\rm nat}-\lambda I)\widehat z+B_Nu=0,\qquad
 A_{CZ}\widehat z+(D_N^{C}-\lambda I)u=0 .
 \tag{NCRP17}
\]
These equations require no inverse of $D_N^C-\lambda I$;
they hold on every original rank stratum.
The original FCI39 residual is exactly the matrix $R_N$ in
NCRP16. In particular
\[
 R_N^*G_NR_N=A_{CZ}^*\Sigma A_{CZ},\qquad
 \operatorname{rank}R_N=\operatorname{rank}A_{CZ}
                         =\operatorname{rank}R_{\mathcal T}>0 .
 \tag{NCRP18}
\]
The last equality follows by taking the quotient of the two full
coefficient frames: $F_\perp^{\,0}$ identifies its coefficient space
with $E/\operatorname{im}L_a$, while the omitted residue-dual
coordinates give another isomorphism of that same quotient.
Their transition is invertible, so both represent the same action
map into that quotient and have the same rank.

Define
\[
 \begin{split}
 D_Z&=A_{\rm nat}^*\Gamma+\Gamma A_{\rm nat}-a\Gamma,\\
 D_C&=(D_N^{C})^*\Sigma+\Sigma D_N^{C}-a\Sigma,\\
 X_N&=\Gamma B_N+A_{CZ}^*\Sigma.
 \end{split}                                                     \tag{NCRP19}
\]
The complete centered same-class equation becomes
\[
 \boxed{\quad
 \widehat z^*D_Z\widehat z+
    2\Re(\widehat z^*X_Nu)+u^*D_Cu
 =2a\delta(\widehat z^*\Gamma\widehat z+u^*\Sigma u).
 \quad}                                                         \tag{NCRP20}
\]
It follows by multiplying the two full blocks in NCRP16 against
the block diagonal metric $\operatorname{diag}(\Gamma,\Sigma)$.
Both leakage directions, $\Gamma B_N$ and $A_{CZ}^*\Sigma$,
remain separately in $X_N$. No vanishing or sign is inferred from
the positivity of $\Gamma$ or $\Sigma$.

The original native coefficient metric is still
$D_{a,N-v}^{(s)}$, at centre $a/2-4$ and order $s=1,a$.
Its exact comparison with the arithmetic section is the positive
matrix
\[
 (D_{a,N-v}^{(s)})^{-1/2}\Gamma
                         (D_{a,N-v}^{(s)})^{-1/2}.
 \tag{NCRP21}
\]
Equations NCRP13--20 do not assign this comparison an isometry or
a bounded condition number. The full conductor square in DNE and
XIM remains the actual morphism connecting those two presentations.

\subsection{The consecutive chart and the first physical metric}

The consecutive chart is used only through the proved chart
transition in the preceding section. In that chart
$\operatorname{im}L_a=\mathbb C[S]_{\le m-1}$.
At $N=q-1$ the remainder quotient has no relation to minimize,
so $G_{q-1}=H_{q-1}^{\rm ar}$ and $m<q$ implies that $Sf$
is still the literal polynomial for each $f$ in this section.
For two such polynomials, the physical identity
$S+\bar S=a$ gives
\[
 \langle f,Sg\rangle_{\rm ar}+
 \langle Sf,g\rangle_{\rm ar}
       =a\langle f,g\rangle_{\rm ar}.
 \tag{NCRP22}
\]
Projecting $Sf,Sg$ onto the same section changes neither of these
pairings. Therefore $D_Z=0$ at this cutoff in this chart.
The full receiver is consequently the exact evaluated identity
\[
 2\Re(\widehat z^*X_{q-1}u)+u^*D_Cu
       =2a\delta\,\|v_\lambda\|_{G_{q-1}}^2>0.
 \tag{NCRP23}
\]
It identifies the entire centered action of the retained class with
the complement and both cross terms. A centered compression by itself
does not bound these terms.

For its residual, let $p_n(y)$ be the monic real orthogonal
polynomials of the original arithmetic measure, with actual norms
$\omega_n=\int_{\mathbb R}|p_n(y)|^2dm_a^{\rm ar}(y)>0$.
Set $\phi_n(S)=p_n((S-c)/i)/\sqrt{\omega_n}$.
Multiplication by the real variable $y$ in this orthonormal frame is
Hermitian. Its coefficient from $\phi_{m-1}$ to $\phi_m$ is
$\sqrt{\omega_m/\omega_{m-1}}$, by comparing the leading coefficients.
For $j<m-1$, $y\phi_j$ has degree below $m$ and has no omitted
component. This proves the full formula
\[
 R_{q-1}
  =i\sqrt{\omega_m/\omega_{m-1}}\,
                       \phi_m e_{m-1}^*,\qquad
 \|R_{q-1}\|=\sqrt{\omega_m/\omega_{m-1}} .
 \tag{NCRP24}
\]
This is the original orthogonal-polynomial recurrence coefficient,
whose conventional formulation is DLMF18.2(iv)
\cite{human:dlmf18}. The proof above fixes its phase and scale.
For the two Gamma measures, and only for them,
$\omega_{n,s}=(2\pi)^{s/2}n!(s/2)_n$, giving
\[
 \|R_{q-1}^{(s)}\|=\sqrt{m(m-1+s/2)},\quad s=1,a,\qquad
 \frac{\|R_{q-1}^{(1)}\|}{a}\to16,\quad
 \frac{\|R_{q-1}^{(a)}\|}{a}\to\sqrt{264}.
 \tag{NCRP25}
\]
The complete arithmetic ratio in NCRP24 has not been replaced by
either Gamma ratio.

For the original arithmetic metric and nonzero
$f\in\mathbb C[S]_{\le m-1}$,
$(A-\lambda)f=(S-\lambda)f$ is literal. Thus
\[
 \begin{split}
 \|(A-\lambda)f\|_{\rm ar}^2
 &=(a\delta)^2\|f\|_{\rm ar}^2\\
 &\quad+\int_{\mathbb R}(y-a\gamma)^2
                      |f(c+iy)|^2dm_a^{\rm ar}(y)
 >(a\delta)^2\|f\|_{\rm ar}^2 .
 \end{split}                                                     \tag{NCRP26}
\]
The strict inequality follows because the integrand can vanish
almost everywhere only if the nonzero polynomial $f$ does.
This is a statement in the consecutive comparison chart at the
first cutoff; no estimate is transported to another chart by
dropping its conductor transition.

\subsection{The original observation projection on the completed class}

Let $\Lambda:E\to B_{\rm obs}$ be any one of the original packet
observation maps, and let $I_K$ be its original full kernel frame.
Its arithmetic orthogonal projector and complement are
\[
 P_K=I_K(I_K^*G_NI_K)^{-1}I_K^*G_N,\qquad P_B=I-P_K.
 \tag{NCRP27}
\]
For a zero-dimensional kernel the first map is zero. For the full
kernel it is the identity. In all cases
$P_K^2=P_K$, $P_K^*G_N=G_NP_K$, and
$P_Kv+P_Bv=v$.
The literal projected vector is
\[
 P_Kv_\lambda
  =P_KF_0u_0+P_KF_C\eta+P_KL_az.
 \tag{NCRP28}
\]
This equation, with the same original $P_K$, is the receiving
correction for any subsequent kernel response. It never replaces
$\ker\Lambda$ by $\operatorname{im}L_a$.

Put $x_K=P_Kv_\lambda$, $x_B=P_Bv_\lambda$ and
$D=A^*G_N+G_NA-aG_N$. Exact substitution of
$Av_\lambda=\lambda v_\lambda$ gives
\[
 \begin{split}
 x_K^*Dx_K
   &=2a\delta\|x_K\|_{G_N}^2
                    -2\Re(x_K^*G_NAx_B),\\
 x_B^*Dx_B
   &=2a\delta\|x_B\|_{G_N}^2
                    -2\Re(x_B^*G_NAx_K),\\
 2\Re(x_K^*Dx_B)
   &=2\Re\{x_K^*G_NAx_B+x_B^*G_NAx_K\}.
 \end{split}                                                     \tag{NCRP29}
\]
For example, $Ax_K=\lambda(x_K+x_B)-Ax_B$ and
$x_K^*G_Nx_B=0$ prove the first line. The second follows by
the other decomposition; in the last line expand $D$ and use
orthogonality to remove only its actual $aG_N$ term.
Adding all three equations returns
$2a\delta\|v_\lambda\|_{G_N}^2$ exactly.
Thus both projected signs still depend on their actual leakage
pairings, now evaluated on the completed class NCRP28.

\subsection{The literal extension of the relative comparison cone}

For four copies let $\mathcal H=\mathbb C^{4m}$ and let
$L=I_4\otimes L_a$. The original coefficient cone is
$B\to W\oplus S\to\mathcal H$ with differentials
$b\mapsto(b,-b)$ and $(w,s)\mapsto w+s$.
The enlarged arithmetic packet cone is obtained by replacing its
last term with $E^{\oplus4}$ and using $L(w+s)$ there.
Since $L$ is injective, its middle cohomology remains
$(S\cap W)/B$. If $T=S+W$, its final cohomology is
$E^{\oplus4}/L(T)$, and the exact sequence is
\[
 0\longrightarrow L(\mathcal H)/L(T)
 \longrightarrow E^{\oplus4}/L(T)
 \longrightarrow E^{\oplus4}/L(\mathcal H)
 \longrightarrow0.
 \tag{NCRP30}
\]
Each map is induced by the identity on representatives. Its exactness
follows by testing whether a representative lies in $L(\mathcal H)$.
The full complement from NCRP3 provides explicit representatives for
the last quotient, of dimension $4(v+q')$.
All quotient metrics are the attained quotients of the full original
$G_N^{\oplus4}$, so their Schur blocks retain the cross terms in
NCRP10 and NCRP13. A retained eigenclass in one indicated copy lies
outside $L(\mathcal H)$ and therefore has a nonzero class in this
enlarged final cohomology.

Multiplication by $S$ on the full packet is the exact action
NCRP9/NCRP16. The native subspace has the nonzero defect NCRP18,
so the action is carried by the complete ambient block matrices.
No unproved preservation of $B,S,W,T$ is used to create an action
on their quotients. In particular NCRP30 is an explicit extension
of the old relative cohomology, with its additional support and
metric; it does not identify a relative denominator with a theta
boundary. The original native determinant allocations remain valid
on their own maps and metrics, and their arithmetic use proceeds
through the completed class and the complete receiving equations above.
